% Bagian: turing-machines
% Bab: machines-computations
% Subbagian: introduction

\documentclass[../../../include/open-logic-section]{subfiles}

\begin{document}

\olfileid[id]{tur}{mac}{int}
\olsection{Pendahuluan}

Apa artinya suatu fungsi, misalnya dari $\Nat$ ke $\Nat$, bersifat
\emph{dapat dihitung}? Salah satu jawaban pertama, sekaligus yang paling
terkenal, ialah bahwa suatu fungsi dapat dihitung jika fungsi itu dapat
dihitung oleh mesin Turing. Gagasan ini dikemukakan oleh Alan Turing pada
1936. Mesin Turing merupakan salah satu contoh \emph{model komputasi}---cara
yang tepat secara matematis untuk mendefinisikan gagasan ``prosedur
komputasi.'' Makna tepat ungkapan itu masih diperdebatkan, tetapi mesin Turing
secara luas disepakati sebagai salah satu cara menentukan prosedur komputasi.
Walaupun istilah ``mesin Turing'' membangkitkan gambaran sebuah mesin fisik
dengan bagian-bagian bergerak, secara ketat mesin Turing merupakan konstruksi
matematis murni dan, karena itu, mengidealkan gagasan prosedur komputasi.
Misalnya, kita tidak membatasi kebutuhan waktu ataupun memori suatu mesin
Turing: mesin Turing dapat menghitung sesuatu sekalipun komputasinya memerlukan
ruang penyimpanan atau jumlah langkah yang melebihi jumlah atom di alam
semesta.

\begin{explain}
Barangkali cara terbaik memandang mesin Turing ialah sebagai program bagi
suatu mekanisme khayalan yang khusus. Mekanisme ini terdiri atas sebuah
\emph{pita} dan sebuah \emph{kepala baca-tulis}. Dalam versi mesin Turing yang
kita gunakan, pita itu tak hingga ke satu arah (ke kanan), dan terbagi menjadi
\emph{petak-petak}; setiap petak dapat memuat sebuah simbol dari suatu
\emph{alfabet} berhingga. Alfabet semacam itu dapat memuat berapa pun banyaknya
simbol yang berlainan, tetapi pada umumnya kita cukup menggunakan tiga simbol:
$\TMendtape$, $\TMblank$, dan $\TMstroke$. Ketika mekanisme dijalankan, pita
berada dalam keadaan kosong (yakni, setiap petak memuat simbol $\TMblank$),
kecuali petak paling kiri yang memuat $\TMendtape$ dan sejumlah berhingga petak
yang memuat \emph{masukan}. Pada setiap saat, mekanisme berada dalam salah satu
dari sejumlah berhingga \emph{keadaan}. Pada awalnya, kepala memindai petak
paling kiri dan mekanisme berada dalam suatu \emph{keadaan awal} yang telah
ditentukan. Pada setiap langkah operasi mekanisme, isi petak yang sedang
dipindai, bersama dengan keadaan mekanisme dan program mesin Turing,
menentukan apa yang terjadi selanjutnya. Program mesin Turing diberikan oleh
suatu fungsi parsial yang menerima keadaan~$q$ dan simbol~$\sigma$ sebagai
masukan serta menghasilkan tripel~$\tuple{q', \sigma', D}$. Setiap kali
mekanisme berada dalam keadaan~$q$ dan membaca simbol~$\sigma$, mekanisme
mengganti simbol pada petak yang sedang dipindai dengan~$\sigma'$; kepala
bergerak ke kiri, ke kanan, atau tetap di tempat menurut apakah $D$ bernilai
$\TMleft$, $\TMright$, atau $\TMstay$; dan mekanisme berpindah ke keadaan~$q'$.

Sebagai contoh, perhatikan situasi pada \olref{fig:tm}.
\begin{figure}
  \olasset{\olpath/assets/diagrams/turing-machine.tikz}
  \caption{Mesin Turing sedang menjalankan programnya.}
  \ollabel{fig:tm}
\end{figure}
Bagian pita mesin Turing yang tampak memuat simbol ujung pita $\TMendtape$ pada
petak paling kiri, diikuti tiga angka~$1$, sebuah angka~$0$, dan empat angka~$1$
lagi. Kepala sedang membaca petak ketiga dari kiri, yang memuat~$1$, dan berada
dalam keadaan~$q_1$---kita mengatakan ``mesin membaca $1$ dalam keadaan~$q_1$.''
Jika, untuk masukan $\tuple{q_1, 1}$, program mesin Turing menghasilkan tripel
$\tuple{q_2, 0, \TMstay}$, mesin itu selanjutnya mengganti~$1$ pada petak ketiga
dengan~$0$, membiarkan kepala baca-tulis tetap di tempatnya, dan berpindah ke
keadaan~$q_2$. Jika kemudian program menghasilkan $\tuple{q_3, 0, \TMright}$
untuk masukan $\tuple{q_2, 0}$, mesin menimpa~$0$ dengan~$0$ yang lain (sehingga
secara efektif membiarkan isi pita di bawah kepala baca-tulis tidak berubah),
bergerak satu petak ke kanan, dan memasuki keadaan~$q_3$. Demikian seterusnya.

Kita mengatakan bahwa mesin \emph{berhenti} ketika menjumpai suatu keadaan
$q_n$ dan simbol~$\sigma$ yang tidak mempunyai instruksi bagi
$\tuple{q_n, \sigma}$, yakni ketika fungsi transisi tidak terdefinisi pada
masukan $\tuple{q_n,\sigma}$. Dengan kata lain, mesin tidak mempunyai instruksi
untuk dijalankan dan pada saat itu berhenti beroperasi. Penghentian kadang-kadang
direpresentasikan oleh suatu keadaan henti khusus~$h$. Hal ini akan diperagakan
secara lebih terperinci nanti.
\end{explain}

\begin{digress}
Keindahan makalah Turing, ``On computable numbers,'' terletak pada kenyataan
bahwa ia tidak hanya menyajikan definisi formal, tetapi juga argumen bahwa
definisi tersebut menangkap gagasan intuitif tentang komputabilitas. Dari
definisi itu, semestinya jelas bahwa setiap fungsi yang dapat dihitung oleh
mesin Turing dapat dihitung dalam arti intuitif. Turing menawarkan tiga jenis
argumen untuk menunjukkan bahwa kebalikannya benar, yakni bahwa setiap fungsi
yang secara wajar kita anggap dapat dihitung dapat dihitung oleh mesin semacam
itu. Ketiganya (menurut kata-kata Turing) adalah:
\begin{enumerate}
\item Seruan langsung kepada intuisi.
\item Bukti kesetaraan dua definisi (apabila definisi baru lebih menarik secara
  intuitif).
\item Pemberian contoh kelas-kelas besar bilangan yang dapat dihitung.
\end{enumerate}
Tujuan kita ialah mencoba mendefinisikan gagasan komputabilitas ``secara
prinsip,'' yakni tanpa memperhitungkan keterbatasan praktis ruang dan waktu.
Tentu saja, setelah definisi komputabilitas yang paling luas tersedia, kita
dapat melanjutkan dengan membahas komputasi yang sumber dayanya dibatasi; hal
ini merupakan inti bidang yang dikenal sebagai ``kompleksitas komputasional.''
\end{digress}

\begin{history}
Alan Turing menemukan mesin Turing pada 1936. Walaupun minatnya ketika itu
adalah keterputusan logika orde pertama, makalah tersebut telah digambarkan
sebagai makalah yang menentukan mengenai landasan perancangan komputer. Dalam
makalah itu, Turing berfokus pada bilangan real yang dapat dihitung, yakni
bilangan real yang ekspansi desimalnya dapat dihitung; tetapi ia mencatat bahwa
gagasannya tidak sulit disesuaikan untuk fungsi dapat dihitung pada bilangan
asli, dan seterusnya. Perhatikan bahwa hal ini terjadi lima tahun penuh sebelum
komputer serbaguna pertama yang berfungsi dibangun pada 1941 (oleh Konrad Zuse
dari Jerman di ruang duduk rumah orang tuanya), tujuh tahun sebelum Turing dan
rekan-rekannya di Bletchley Park membangun mesin pemecah sandi Colossus (1943),
sembilan tahun sebelum ENIAC buatan Amerika Serikat (1945), dua belas tahun
sebelum komputer serbaguna Inggris pertama---Manchester Small-Scale
Experimental Machine---dibangun di Manchester (1948), dan tiga belas tahun
sebelum orang Amerika pertama kali menguji BINAC (1949). Manchester SSEM
menyandang kehormatan sebagai komputer pertama yang menyimpan program---mesin
sebelumnya harus dirangkai ulang dengan tangan untuk setiap tugas baru.
\end{history}

\end{document}
